% Part: first-order-logic
% Chapter: natural-deduction
% Section: identity

\documentclass[../../../include/open-logic-section]{subfiles}

\begin{document}

\olfileid[fr]{fol}{ntd}{ide}

\olsection{\usetoken{P}{derivation} avec égalité}

Les !!{derivation}s avec !!{identity} nécessitent des règles
d'inférence supplémentaires.
\begin{defish}
\AxiomC{}
\RightLabel{\Intro{\eq}}
\UnaryInfC{$\eq[t][t]$}
\DisplayProof
\hfill
\begin{tabular}{r}
\AxiomC{$\eq[t_1][t_2]$}
\AxiomC{$!A(t_1)$}
\RightLabel{\Elim{\eq}}
\BinaryInfC{$!A(t_2)$}
\DisplayProof
\\[3ex]
\AxiomC{$\eq[t_1][t_2]$}
\AxiomC{$!A(t_2)$}
\RightLabel{\Elim{\eq}}
\BinaryInfC{$!A(t_1)$}
\DisplayProof
\end{tabular}
\end{defish}
Dans ces règles, $t$, $t_1$ et $t_2$ sont des termes clos.
La règle \Intro{\eq} permet de !!{derive} directement tout énoncé
d'égalité de la forme $\eq[t][t]$, sans aucune hypothèse.

\begin{ex}
Si $s$ et $t$ sont des termes clos, alors
$!A(s), \eq[s][t] \Proves !A(t)$ :
\begin{prooftree}
\AxiomC{$\eq[s][t]$}
\AxiomC{$!A(s)$}
\RightLabel{$\Elim{\eq}$}
\BinaryInfC{$!A(t)$}
\end{prooftree}
On connaît aussi ce principe sous le nom de « substituabilité
des égaux », ou loi de Leibniz.
\end{ex}

\begin{prob}
Démontrer que $=$ est symétrique et transitif : donner des
!!{derivation}s de $\lforall[x][\lforall[y][(\eq[x][y] \lif
    \eq[y][x])]]$ et de $\lforall[x][\lforall[y][\lforall[z][((\eq[x][y]
    \land \eq[y][z]) \lif \eq[x][z])]]]$.
\end{prob}

\begin{ex}
Nous allons !!{derive} l'!!{sentence}
\begin{align*}
& \lforall[x][\lforall[y][((!A(x) \land !A(y)) \lif \eq[x][y])]]
\intertext{à partir de l'!!{sentence}}
& \lexists[x][\lforall[y][(!A(y) \lif \eq[y][x])]]
\end{align*}
Choisissons des constantes $a$, $b$ et $c$ deux à deux distinctes
qui ne figurent pas dans cette hypothèse, puis construisons la
!!{derivation} en remontant depuis sa conclusion :
\begin{prooftree}
\AxiomC{$\lexists[x][\lforall[y][(!A(y) \lif \eq[y][x])]]
\quad \Discharge{!A(a) \land !A(b)}{1}$}
\DeduceC{$\eq[a][b]$}
\DischargeRule{\Intro{\lif}}{1}
\UnaryInfC{$((!A(a) \land !A(b)) \lif \eq[a][b])$}
\RightLabel{\Intro{\lforall}}
\UnaryInfC{$\lforall[y][((!A(a) \land !A(y)) \lif \eq[a][y])]$}
\RightLabel{\Intro{\lforall}}
\UnaryInfC{$\lforall[x][\lforall[y][((!A(x) \land !A(y)) \lif \eq[x][y])]]$}
\end{prooftree}
Il faut maintenant utiliser l'hypothèse principale. Comme c'est
une !!{formula} existentielle, nous appliquons \Elim{\lexists}
pour !!{derive} la conclusion intermédiaire $\eq[a][b]$.
\begin{prooftree}
\AxiomC{$\lexists[x][\lforall[y][(!A(y) \lif \eq[y][x])]]$}
\AxiomC{$\Discharge{\lforall[y][(!A(y) \lif \eq[y][c])]}{2}$}
\noLine
\UnaryInfC{$\Discharge{!A(a) \land !A(b)}{1}$}
\DeduceC{$\eq[a][b]$}
\DischargeRule{\Elim{\lexists}}{2}
\BinaryInfC{$\eq[a][b]$}
\DischargeRule{\Intro{\lif}}{1}
\UnaryInfC{$((!A(a) \land !A(b)) \lif \eq[a][b])$}
\RightLabel{\Intro{\lforall}}
\UnaryInfC{$\lforall[y][((!A(a) \land !A(y)) \lif \eq[a][y])]$}
\RightLabel{\Intro{\lforall}}
\UnaryInfC{$\lforall[x][\lforall[y][((!A(x) \land !A(y)) \lif \eq[x][y])]]$}
\end{prooftree}
Pour compléter la sous-!!{derivation} en haut à droite, on
utilise ses hypothèses afin d'établir $\eq[a][c]$ et $\eq[b][c]$.
Il faut deux !!{derivation}s distinctes. Celle de $\eq[a][c]$
est la suivante :
\begin{prooftree}
\AxiomC{$\Discharge{\lforall[y][(!A(y) \lif \eq[y][c])]}{2}$}
\RightLabel{\Elim{\lforall}}
\UnaryInfC{$!A(a) \lif \eq[a][c]$}
\AxiomC{$\Discharge{!A(a) \land !A(b)}{1}$}
\RightLabel{\Elim{\land}}
\UnaryInfC{$!A(a)$}
\RightLabel{\Elim{\lif}}
\BinaryInfC{$\eq[a][c]$}
\end{prooftree}
À partir de $\eq[a][c]$ et $\eq[b][c]$, on !!{derive}
$\eq[a][b]$ par \Elim{\eq}.
\end{ex}

\begin{prob}
Donner des !!{derivation}s des !!{formula}s suivantes :
\begin{enumerate}
\item $\lforall[x][\lforall[y][((\eq[x][y] \land !A(x)) \lif !A(y))]]$
\item $\lexists[x][!A(x)] \land \lforall[y][\lforall[z][((!A(y) \land
    !A(z)) \lif \eq[y][z])]] \lif \lexists[x][(!A(x) \land
  \lforall[y][(!A(y) \lif \eq[y][x])])]$
\end{enumerate}
\end{prob}

\begin{editorial}
Dans l'exercice sur la transitivité, le troisième quantificateur
de la source n'encadre pas correctement son argument; les crochets
sont rétablis. Dans deux arbres du deuxième exemple, la parenthèse
de l'hypothèse universelle est replacée à l'intérieur de l'argument
du quantificateur. Le choix de trois constantes fraîches et
distinctes, nécessaire aux conditions sur les eigenvariables,
est explicité. Les deux orientations de l'élimination de l'égalité
sont conservées.
\end{editorial}

\end{document}
